% s-347-p347-m002; original PDF 347, 348, 349
\Aphorism{s-347-p347-m002}{399}
Schopenhauer, der nicht nur als Farbenlehrer, sondern zugleich als Hasser der Mathematik auf Goethes Spuren wandelt, sagt geradezu, dass nur das Was und Wie wissenschaftlich in Betracht komme — das Wieviel und Wiegross gehe allein die Praxis an und könne im Übrigen mit einer „ungefähren Schätzung“ abgethan werden. Der Fuchs und die Trauben! Mit der quantitativen Forschung beginnt erst Strenge, Baugrund, Methode, aber auch Kampf und Schwierigkeit, während nichts leichter gebaut ist und leichter umfällt als eine bloss qualitative Hypothese. Der „königliche Weg“ zu den Dingen ist kürzer als der analytische Dornenpfad, aber dieser ist „schwindel“freier; in nichts kann man sich eclatanter vergreifen als in einer oberflächlichen Schätzung des Was und Wie, ohne genaue Berechnung des Wieviel. Um ein Beispiel, ein freilich etwas enges Beispiel zu nennen: wenn es im Weltraum ein widerstehendes Mittel (Äther) giebt, so wird dadurch die Tangentialgeschwindigkeit eines Kometen vermindert, also an sich seine Umlaufszeit \Sperr{vergrössert} — andererseits gewinnt dadurch die Anziehung der Sonne an Übergewicht, bringt den Kometen dem Centrum näher und \Sperr{verkleinert} seine Umlaufszeit. Welcher dieser beiden Effecte behält die Oberhand? Ich fordere jeden \Greek{ἀγεωμέτρητος} auf, durch rein qualitative Abschätzung festzustellen, ob die Umlaufszeit unter dem Einfluss der beiden entgegengesetzten Antriebe thatsächlich grösser oder kleiner wird! — Machen wir uns nichts vor, wir Erkennenden! Es giebt kein qualitatives Wissen in Dingen, die der Messung und Rechnung zugänglich sind; das Was und Wie sagt gar nichts, solange man nicht auch das Wieviel kennt, und zwar genau und gründlich kennt. Jede qualitative Erklärung von Thatbeständen kommt auf den Punkt, wo sie gegen andere, ebenso plausible Erklärungen die quantitative Probe bestehen muss, nämlich die Probe auf \Sperr{Congruenz} \Sperr{mit} \Sperr{der} \Sperr{Erfahrung} bis zum Hundertel Secunde und Milliontel Millimeter hinunter. Auf Grund eines solchen messenden und rechnenden \Foreign{experimentum crucis} siegte auch über Newtons Farbenlehre — nicht die Goethische oder Schopenhauerische, sondern die Huygens’sche Vibrationstheorie. Gegen diese, die ihm völlig Idiosynkrasie und unverständlich ist, protestirt Schopenhauer leidenschaftlich und mit derbem Hohn über die Zumuthung, sich den ruhigen, selbst im Sturme unbewegten Lichtstrahl als „Trommelwirbel der Äthertremulanten“ vorzustellen — man sieht, die bloss „qualitative“ Naturbetrachtung mündet nicht weit vom „gesunden Menschenverstand“, der sich auch in seiner Weise das Wie und Was zurechtlegt und über das Wieviel keine grauen Haare wachsen lässt. Das ist die Vernünftelei über das Wirkliche, das Spintisiren und Meditiren ohne Beobachtung, das schon Aristophanes in den „Wolken“ verspottet; diese Anthropomorphismen und Plausibilitäts-Erwägungen haben die Wissenschaft länger aufgehalten als Inquisition und russische Censur. Heute wird uns kein Laienvorurtheil mehr hindern, Ruhe durch Bewegung, geradlinige Kräfte durch Undulation, Einfaches durch Complicirtes zu erklären, sobald damit ein besserer Anschluss an die beobachtete, zahlenmässig festgelegte Wirklichkeit erreicht wird. Gewiss ist ein geradlinig fortgeschleuderter Lichtstoff dem populären Gutdünken einleuchtender als eine sich fortpflanzende Welle in der noch nie gesehenen imaginären Flüssigkeit „Äther“; dennoch ist er zu verwerfen, sobald er den wirklichen Thatbestand nicht genau genug darstellt. Wir werden uns nicht bedenken, an Stelle jeder „plausiblen“ und „natürlichen“ Hypothese eine höchst verzwickte Wellen-, Strudel-, Wirbelhypothese zu setzen, sobald sie \Sperr{metrisch} mehr leistet als die alte; wir werden das Was und Wie, so wie es sich im \Foreign{common sense} malt, schlankweg auf den Kopf stellen, wenn wir dadurch irgend einen Widerspruch von 0,003 auf 0,002 herabdrücken können.
\par

